% shared/preamble.tex
% -----------------------------------------------------------------------------
% Single source of truth for packages, theorem environments, macros, and the
% bibliography resource. Inputted by main.tex via
%   \usepackage{subfiles}
%   \input{\subfix{shared/preamble.tex}}
% so that `latexmk main.tex`, a standalone `latexmk modules/NN-*.tex`, and a
% standalone `latexmk solutions/<id>.tex` all resolve the same preamble and
% bibliography. Do NOT put \documentclass or \usepackage{subfiles} here (those
% live in main.tex).
%
% Build (biber backend), from the repository root:
%   latexmk -pdf -outdir=build main.tex          % the whole document
% A standalone module shows cross-module \ref as "??" and unresolved citations;
% this is expected and harmless for syntax/math checking.
%
% Macros are in two blocks: a domain-neutral core, then this repository's own
% program macros. Add yours to the second block so the first stays diffable
% against the template.
% -----------------------------------------------------------------------------

\usepackage[margin=1in]{geometry}
\usepackage[T1]{fontenc}
\usepackage[utf8]{inputenc}
\usepackage{lmodern}
\usepackage{microtype}
\usepackage{amsmath,amssymb,amsthm,mathtools}
\usepackage{mathrsfs}
\usepackage{bm}
\usepackage{enumitem}
\usepackage{booktabs}
\usepackage{tabularx}
\usepackage{array}
\usepackage{xcolor}
\usepackage{tikz}
\usetikzlibrary{arrows.meta,positioning,fit,backgrounds}
\usepackage{hyperref}
\usepackage[backend=biber,style=alphabetic,maxbibnames=8]{biblatex}
\addbibresource{\subfix{references.bib}}

\hypersetup{
  colorlinks=true,
  linkcolor=blue!55!black,
  citecolor=blue!55!black,
  urlcolor=blue!55!black
}

% --- theorem environments (shared numbering by section) ----------------------
\newtheorem{theorem}{Theorem}[section]
\newtheorem{lemma}[theorem]{Lemma}
\newtheorem{proposition}[theorem]{Proposition}
\newtheorem{corollary}[theorem]{Corollary}
\newtheorem{conjecture}[theorem]{Conjecture}
\newtheorem{question}[theorem]{Question}
\newtheorem{claim}[theorem]{Claim}
\theoremstyle{definition}
\newtheorem{definition}[theorem]{Definition}
\newtheorem{assumption}[theorem]{Assumption}
\newtheorem{hypothesis}[theorem]{Hypothesis}
\newtheorem{program}[theorem]{Program}
\newtheorem{family}[theorem]{Family}
\theoremstyle{remark}
\newtheorem{remark}[theorem]{Remark}
\newtheorem{warning}[theorem]{Warning}
\newtheorem{example}[theorem]{Example}
\newtheorem{heuristic}[theorem]{Heuristic}

\newcolumntype{Y}{>{\raggedright\arraybackslash}X}

% --- group dividers ----------------------------------------------------------
% \partgroup{<heading>}{<subtitle>} emits a visible divider plus a
% table-of-contents entry ABOVE \section level, so that a long document's grouping
% of sections is legible both in the TOC and on the page. It introduces no
% numbering of its own: section numbers, theorem numbers, and every \label are
% untouched, which matters because a \label IS a ledger node id (see CLAUDE.md).
\makeatletter
\newcommand{\l@partgroup}[2]{%
  \ifnum \c@tocdepth >-2\relax
    \addpenalty{-\@highpenalty}%
    \vskip 1.0em \@plus\p@
    \begingroup
      \parindent \z@ \rightskip \@pnumwidth
      \parfillskip -\@pnumwidth
      {\leavevmode\bfseries #1}\hfil \nobreak\hb@xt@\@pnumwidth{\hss #2}\par
      \nobreak
    \endgroup
  \fi}
\makeatother
\newcommand{\partgroup}[2]{%
  \par\addvspace{2.0\baselineskip}%
  \phantomsection
  \noindent\textcolor{blue!55!black}{\rule{\textwidth}{1.1pt}}\par\vspace{0.30em}%
  {\noindent\large\bfseries #1\par}%
  \vspace{0.15em}%
  {\noindent\itshape #2\par}%
  \vspace{0.30em}%
  \noindent\textcolor{blue!55!black}{\rule{\textwidth}{0.4pt}}\par
  \addcontentsline{toc}{partgroup}{#1}%
  \addvspace{1.0\baselineskip}%
}

% --- core macros (domain-neutral) --------------------------------------------
\newcommand{\R}{\mathbb{R}}
\newcommand{\E}{\mathbb{E}}
\newcommand{\Prob}{\mathbb{P}}
\newcommand{\one}{\mathbf{1}}
\newcommand{\Cov}{\operatorname{Cov}}
\newcommand{\Var}{\operatorname{Var}}
\newcommand{\Tr}{\operatorname{Tr}}
\newcommand{\Ent}{\operatorname{Ent}}
\newcommand{\KL}{\operatorname{KL}}
\newcommand{\Id}{\mathrm{Id}}
\newcommand{\op}{\mathrm{op}}
\newcommand{\osc}{\operatorname{osc}}
\newcommand{\diag}{\operatorname{diag}}
\newcommand{\Hess}{\nabla^2}
\newcommand{\dd}{\,\mathrm{d}}
\newcommand{\eps}{\varepsilon}
\newcommand{\lmax}{\lambda_{\max}}
\newcommand{\lmin}{\lambda_{\min}}

% --- program macros ----------------------------------------------------------
% This repository's own notation. Everything above this line is domain-neutral
% and should stay diffable against the template; everything here is yours.
% Record the fixed normalizations these macros stand for in shared/notation.md.


% --- delimiters --------------------------------------------------------------
\newcommand{\norm}[1]{\left\lVert #1\right\rVert}
\newcommand{\abs}[1]{\left\lvert #1\right\rvert}
\newcommand{\inner}[2]{\left\langle #1,#2\right\rangle}
